% secciones/03_diseno_robot.tex
\section{Diseño y Características del Robot}
El diseño de un robot destinado al entorno hospitalario debe concebirse desde una perspectiva holística que equilibre la funcionalidad logística con la seguridad biológica y la aceptabilidad social. A diferencia de los robots industriales convencionales, un dispositivo de asistencia sanitaria debe operar en un entorno semi-estructurado y altamente dinámico, lo que impone requisitos mecánicos y electrónicos extremadamente rigurosos.

En lo que respecta a la estructura externa, la elección de materiales es crítica para garantizar la asepsia. La carcasa del robot se proyecta utilizando polímeros de grado médico con propiedades antimicrobianas intrínsecas o acero inoxidable de alta calidad, permitiendo una desinfección profunda mediante agentes químicos agresivos como el peróxido de hidrógeno vaporizado o el alcohol isopropílico. El diseño industrial prioriza las superficies lisas y los acabados redondeados, una decisión técnica orientada a eliminar puntos de acumulación de patógenos y a minimizar el impacto mecánico en caso de una colisión accidental con el mobiliario hospitalario o el personal. Para la logística de materiales sensibles, el robot integra compartimentos modulares sellados con sistemas de cierre electromecánico, cuya apertura se gestiona mediante autenticación segura por RFID o biometría, asegurando la trazabilidad total de los medicamentos o muestras transportadas.

La arquitectura interna de sensores constituye el sistema nervioso del robot, permitiendo una percepción precisa del entorno para la navegación autónoma. El sistema principal de localización y mapeo simultáneo (SLAM) se apoya en una unidad LiDAR 2D de alta frecuencia que genera un mapa planar del entorno, complementada por cámaras de profundidad RGB-D. Estas cámaras son fundamentales para detectar obstáculos que escapan al plano del LiDAR, tales como pies de suero, camillas o extremidades de pacientes en reposo. Como medida de redundancia y seguridad ante fallos, se incorporan sensores de ultrasonido perimetrales que actúan como un "parachoques virtual", deteniendo el movimiento de forma inmediata ante una proximidad crítica.

Uno de los elementos diferenciadores de nuestra propuesta es el sistema de locomoción y control, fundamentado en una configuración de cuatro ruedas Mecanum. Esta elección técnica otorga al robot omnidireccionalidad total, permitiendo desplazamientos en los tres grados de libertad del plano ($x, y, \theta$). Esta capacidad de traslación lateral y rotación sobre el propio eje es crítica para maniobrar en espacios confinados, como el interior de ascensores o pasillos congestionados, facilitando el paso de camillas sin interrumpir la navegación. Cada rueda es accionada de forma independiente por motores síncronos de imanes permanentes o \textit{Brushless} (BLDC) con reductoras planetarias. Esta combinación asegura un funcionamiento extremadamente silencioso —fundamental para el confort acústico de los pacientes— y proporciona un par motor constante capaz de movilizar con agilidad la carga útil de 25 kg y el peso estructural del chasis.

Para proteger las muestras biológicas de las vibraciones y garantizar la precisión cinemática, se ha diseñado una suspensión independiente de tipo \textit{Double Wishbone} (doble horquilla) de grado automotriz. Dado que la eficacia de las ruedas Mecanum depende de que los rodillos mantengan un contacto constante y uniforme con el suelo, este sistema de doble brazo con muelle-amortiguador compensa las irregularidades del pavimento y los resaltes en juntas de dilatación. Esto evita que el robot pierda tracción o se desoriente ante pequeños baches, a la vez que aísla la carga sensible de impactos mecánicos. 

Finalmente, la sensorización de la locomoción aborda el problema del deslizamiento (\textit{slippage}) inherente a los rodillos de las ruedas Mecanum, especialmente en suelos hospitalarios pulidos. La odometría del sistema no depende exclusivamente de los \textit{encoders} de los motores, sino que se articula mediante una arquitectura de fusión sensorial. Se implementa un Filtro de Kalman Extendido (EKF) que integra en tiempo real las lecturas de los \textit{encoders} con los datos de una unidad de medición inercial (IMU) de 9 ejes. Este procesamiento, gestionado de forma nativa en ROS 2, permite corregir las desviaciones acumuladas y estimar la posición real del robot con alta fidelidad, asegurando que la navegación sea fluida, predecible y segura en todo momento.

Para garantizar la versatilidad operativa definida en los requerimientos funcionales, el chasis integra un brazo robótico colaborativo de bajo peso y alta precisión. Este manipulador está estratégicamente situado en la parte superior del módulo motriz, permitiendo un rango de alcance que cubre desde la altura de los botones de un ascensor estándar hasta el pomo de una puerta hospitalaria. El diseño del efector final (pinza) es de tipo adaptativo, lo que le permite interactuar con elementos físicos del entorno sin necesidad de una domotización previa de la infraestructura. Este sistema actúa como una capa de redundancia física: aunque el robot prioriza la comunicación inalámbrica vía Wi-Fi para la gestión de accesos, el brazo asegura que el flujo logístico no se interrumpa en zonas con conectividad limitada o infraestructuras antiguas.

El compartimento de carga, diseñado para una capacidad nominal de 25 kg, constituye el núcleo del servicio del robot. Interiormente, la cabina está revestida con materiales reflectantes y equipada con un sistema de lámparas UV-C de alta intensidad. Siguiendo el requerimiento de asepsia activa, el sistema de control activa un ciclo de desinfección automática tras cada entrega. Este proceso está blindado mediante sensores de cierre electromecánico que impiden la emisión de radiación ultravioleta si la compuerta no está perfectamente sellada, garantizando así la seguridad del personal sanitario. El volumen interno está optimizado para el transporte de gradillas de laboratorio y contenedores de medicación, manteniendo la carga en una posición central para no desplazar el centro de gravedad del conjunto y comprometer la estabilidad durante el frenado.

En cuanto a la seguridad activa, el diseño contempla un sistema de frenado electromecánico redundante. Para cumplir con la restricción de una distancia de detención inferior a 50 cm a velocidad máxima, el controlador de movimiento implementa una curva de desaceleración de emergencia que actúa simultáneamente sobre los inversores de los motores BLDC y un freno físico de seguridad. Esta capacidad de respuesta se complementa con el algoritmo de "Navegación Social", el cual redefine el espacio personal del robot. Mediante el uso de zonas de coste dinámicas en el planificador de rutas de ROS 2, el sistema garantiza un margen de respeto de 30 cm respecto a cualquier ser humano detectado por las cámaras de profundidad, asegurando que el despliegue del robot sea percibido como una presencia segura y no intrusiva en las plantas de hospitalización.

% --- Añadir en la Sección 3, justo antes del apartado de Dimensionamiento Energético ---

Para orquestar la compleja arquitectura de hardware descrita y garantizar el cumplimiento de la computación \textit{Edge} en tiempo real, el cerebro del robot se fundamenta en un módulo de procesamiento integrado \textbf{NVIDIA Jetson Orin NX}. La elección de esta arquitectura híbrida (CPU ARM + GPU Ampere) frente a un procesador x86 industrial tradicional responde a la necesidad de equilibrar una alta capacidad de cómputo con una alta eficiencia energética.

Este módulo actúa como el núcleo del sistema operativo ROS 2, gestionando los algoritmos críticos de baja latencia sin depender de la nube. Por un lado, su capacidad de aceleración por hardware (mediante el entorno \textit{NVIDIA Isaac ROS}) permite procesar fluidamente las densas nubes de puntos 3D generadas por las cámaras RGB-D y el LiDAR, lo cual es estrictamente necesario para ejecutar las redes neuronales de detección de personas y cumplir con el margen de 30 cm de la navegación social (RNF-7). Por otro lado, la unidad de procesamiento central absorbe la carga computacional de la cinemática inversa del brazo robótico y la fusión sensorial (Filtro de Kalman Extendido) de las ruedas Mecanum. Todo ello operando bajo un TDP (\textit{Thermal Design Power}) máximo de 25 W, lo que minimiza el impacto térmico en el chasis cerrado y optimiza el consumo energético global del dispositivo.

% --- Añadir al final de la Sección 3 ---

\subsection*{Dimensionamiento Energético y Capacidad de Batería}
Para garantizar el cumplimiento del requerimiento RNF-1 (operación ininterrumpida de al menos 8 horas en franjas de alta demanda) y soportar la arquitectura de hardware descrita, es imperativo realizar un dimensionamiento energético preliminar. El cálculo se basa en un balance de potencias promedio, asumiendo un peso total del sistema de 75 kg (50 kg de chasis y componentes + 25 kg de carga útil máxima) y un voltaje de operación nominal del bus de 24 V.

El consumo total estimado ($P_{total}$) se divide en tres subsistemas principales:
\begin{itemize}
    \item \textbf{Locomoción ($P_{loc}$):} Para desplazar una masa de 75 kg a una velocidad nominal de 1.2 m/s, asumiendo un ciclo de trabajo del 60\% (el robot se mueve el 60\% del tiempo y está en espera el 40\%), se estima un consumo medio de 120 W.
    \item \textbf{Computación y Sensórica ($P_{sys}$):} La arquitectura \textit{Edge} con un procesador NVIDIA Jetson, junto con el LiDAR, las cámaras de profundidad RGB-D, la IMU y los módulos de comunicación Wi-Fi, presentan un consumo constante y elevado. Se estima un requerimiento continuo de 60 W.
    \item \textbf{Periféricos de Interacción ($P_{per}$):} El brazo robótico y la lámpara de desinfección UV-C tienen un consumo pico alto, pero un ciclo de trabajo muy bajo (inferior al 5\% del turno). Se les asigna una potencia promedio ponderada de 15 W.
\end{itemize}

El consumo de potencia media ($P_{avg}$) del sistema en operación es la suma de estos factores:
\begin{equation}
    P_{avg} = P_{loc} + P_{sys} + P_{per} = 120 \text{ W} + 60 \text{ W} + 15 \text{ W} = 195 \text{ W}
\end{equation}

Para cumplir con el turno de 8 horas ($t = 8$ h), la energía total requerida ($E$) es:
\begin{equation}
    E = P_{avg} \times t = 195 \text{ W} \times 8 \text{ h} = 1560 \text{ Wh}
\end{equation}

Con el fin de maximizar la vida útil de las celdas y evitar que el robot se apague abruptamente, no se debe descargar la batería por debajo del 20\% de su capacidad (Profundidad de Descarga o \textit{DoD} del 80\%). Por tanto, la energía bruta de la batería ($E_{bruta}$) debe ser:
\begin{equation}
    E_{bruta} = \frac{E}{DoD} = \frac{1560 \text{ Wh}}{0.80} = 1950 \text{ Wh}
\end{equation}

Traduciendo esta energía a capacidad eléctrica ($C$) a la tensión nominal del sistema (24 V):
\begin{equation}
    C \text{ (Ah)} = \frac{E_{bruta}}{24 \text{ V}} = 81.25 \text{ Ah}
\end{equation}

\textbf{Conclusión de Diseño:} En base a estos cálculos, el diseño final prescribirá un banco de baterías de tecnología \textbf{LiFePO4 (Litio-Ferrofosfato)} con una capacidad comercial estándar de \textbf{24 V y 100 Ah} (aproximadamente 2.4 kWh). Aunque el requerimiento mínimo calculado es de 81.25 Ah, se opta por el modelo de 100 Ah para proporcionar un margen de seguridad del 18\%, compensando la posible degradación de las celdas tras ciclos de uso intensivos y garantizando que el robot pueda completar misiones críticas de última hora incluso al final de su turno sin comprometer el umbral de seguridad de descarga. Esta química es ideal por su nulo riesgo de incendio en entornos críticos y su soporte para cargas de oportunidad rápidas (RF-9).